\documentclass[11pt]{article}
\usepackage{amsmath,amssymb}
\usepackage{physics}
\usepackage[margin=1in]{geometry}
\usepackage{enumitem}

\title{ATPL 2025 --- Exercises: Classical Simulability and the Limits of
Quantum Advantage}
\author{Lecture 1: Why do we need hybrid programming models?}
\date{}

\begin{document}
\maketitle

Background: Gottesman-Knill theorem; Aaronson \& Gottesman,
\emph{Improved Simulation of Stabilizer Circuits} (2004); Nielsen \&
Chuang, Section 10.5 (the stabilizer formalism).

\begin{enumerate}[label=\textbf{Exercise \arabic*.}, leftmargin=2.2cm]

\item \textbf{Clifford conjugation rules.}
Recall that $H$, $S=R_Z(\pi/2)$, and CNOT map Pauli operators to
(signed) Pauli operators under conjugation, e.g.
\[
HXH = Z, \qquad HZH = X, \qquad SXS^\dagger = Y \ \text{(up to a sign
convention)},
\]
and CNOT (control on qubit 1) maps
\[
X\otimes I \mapsto X\otimes X, \quad
I\otimes X \mapsto I\otimes X, \quad
Z\otimes I \mapsto Z\otimes I, \quad
I\otimes Z \mapsto Z\otimes Z.
\]
\begin{enumerate}[label=(\alph*)]
\item Verify at least two of these relations by direct $2\times2$ (or
$4\times4$) matrix multiplication.
\item Explain why this "Pauli-to-Pauli" property is exactly what lets a
stabilizer simulator track an $n$-qubit state using only $\mathcal O(n)$
Pauli generators (each specified by $\mathcal O(n)$ bits) rather than
$2^n$ complex amplitudes.
\end{enumerate}

\item \textbf{Tracking a Bell-state circuit via stabilizers.}
The state $\ket{00}$ is stabilized by the generators
$\{Z\otimes I,\ I\otimes Z\}$ (each has $\ket{00}$ as a $+1$
eigenstate, and together with a choice of signs they uniquely specify
the state).
\begin{enumerate}[label=(\alph*)]
\item Apply $H$ to qubit 1: using the conjugation rules from Exercise 1,
find the new stabilizer generators of $(H\otimes I)\ket{00}$.
\item Apply CNOT (qubit 1 control, qubit 2 target) to the result: find
the new stabilizer generators, and confirm they are
$\{X\otimes X,\ Z\otimes Z\}$ -- the stabilizers of the Bell state
$\ket{\Phi^+}$.
\item Contrast the amount of information used in this derivation (two
short Pauli strings) with the amount needed to write out the state
vector directly (4 complex amplitudes). For $n=2$ this gap is small --
explain, with a formula, how it grows as a function of $n$ for an
$n$-qubit stabilizer state built the same way.
\end{enumerate}

\item \textbf{Why $T$ breaks the tableau method.}
The $T$ gate is $R_Z(\pi/4)$ (up to global phase).
\begin{enumerate}[label=(\alph*)]
\item Compute $TXT^\dagger$ explicitly as a $2\times2$ matrix.
\item Show that the result is \emph{not} (up to a global phase and
sign) equal to any single Pauli matrix $I,X,Y,Z$. What does this mean
for trying to track the effect of a $T$ gate using only a Pauli-string
tableau?
\end{enumerate}

\item \textbf{Counting non-Clifford gates.}
Using the Gottesman--Aaronson bound quoted in the lecture,
$\mathcal O\!\left(n^2 + nM + nT\,2^{\gamma T}\right)$ time to classically
simulate a circuit with $M$ Clifford gates and $T$ non-Clifford (e.g.\
$T$-gate) operations on $n$ qubits:
\begin{enumerate}[label=(\alph*)]
\item For $n=50$ qubits and $M=10^4$ Clifford gates, describe
qualitatively how the simulation time scales as $T$ grows from $0$ to
$10$, $20$, then $40$. (You don't need $\gamma$'s exact value -- just
discuss the behaviour of the $2^{\gamma T}$ term relative to the
polynomial terms.)
\item Explain, in your own words, why this gives a \emph{necessary but
not sufficient} pair of conditions for quantum advantage (the lecture's
two-part conclusion). Sketch an example of a problem that plausibly
satisfies condition 1 (small classical input/output) but might fail
condition 2 (i.e.\ it can be solved with few non-Clifford gates), or vice
versa.
\end{enumerate}

\item \textbf{BQP, NP, and Grover.}
Grover's algorithm gives a quadratic speedup for unstructured search,
turning $\mathcal O(2^n)$ classical query complexity into
$\mathcal O(2^{n/2})$. Explain briefly why this quadratic speedup is
not believed to be enough to place NP-complete problems inside BQP
(i.e.\ why BQP is conjectured to still exclude NP-hard problems, as in
the "sweet spot" diagram from the lecture), even though unstructured
search is closely related to many NP-complete problems.

\end{enumerate}

\end{document}
